\documentclass[11pt]{article}
\usepackage{notes} % place notes.sty in the same folder
\definecolor{AccN}{HTML}{1B7A5E}
\fancyhead[L]{\small\color{NavyN}\nouppercase{\leftmark}}\fancyhead[R]{\small\color{GreyN}Instructor Manual --- Lecture 7}
\usepackage{longtable,multicol,needspace}
\setlength{\emergencystretch}{2em}
\newtcolorbox{instr}[1]{nb,colback=AccN!6,colframe=AccN,title={#1}}
\newtcolorbox{prob}{nb,colback=BluN!5,colframe=BluN,title={Problem statement}}
\newtcolorbox{soln}{nb,colback=GrnN!5,colframe=GrnN,title={Model solution and marking points}}
\newtcolorbox{sayit}{nb,colback=white,colframe=NavyN,title={Read aloud to students}}
\newtcolorbox{live}[1]{nb,colback=AccN!4,colframe=AccN,title={#1}}
\begin{document}
\begin{titlepage}\centering\vspace*{2.2cm}
{\color{AccN}\large\bfseries LECTURE 7\par}\vspace{0.6cm}
{\color{NavyN}\Huge\bfseries Reinforcement Learning Part 2:\\Policy Gradient Methods\par}\vspace{0.6cm}
{\color{AccN}\rule{0.7\textwidth}{1.5pt}}\par\vspace{0.6cm}
{\LARGE Instructor Manual\par}\vspace{0.4cm}
{\large\begin{minipage}{0.85\textwidth}\centering REINFORCE \textbullet\ baselines \textbullet\ actor--critic \textbullet\ A2C/A3C \textbullet\ PPO \textbullet\ continuous control \textbullet\ SAC\\Case study: 4-drone autonomous swarm, continuous waypoint control, 6-DoF flight dynamics, PPO vs SAC\end{minipage}\par}\vspace{2cm}
{\large AK DHAMIJA, Researcher \& Instructor, AI \& Machine Learning\par}\vspace{0.2cm}
{\normalsize \href{mailto:dhamija.ak@gmail.com}{dhamija.ak@gmail.com} \quad\textbullet\quad 9811453163\par}\vspace{0.3cm}
{Unit 4: Simulation Strategies in Warfare --- AI \& Machine Learning in Warfare\par}\vfill
\begin{tcolorbox}[colback=NavyN!4,colframe=NavyN,boxrule=0.6pt,arc=2pt,width=0.9\textwidth]\small
\textbf{This manual contains:} the 45-minute session plan with timings; a slide-by-slide teaching guide for all 19 slides; key messages, facilitation tips, anticipated questions with model answers, misconceptions and in-class activities; live demonstrations with the interactive studio; consistency notes on the source material; the complete quiz answer key with rationale and scoring rules; and the assignment with marking scheme and full model solutions.
\end{tcolorbox}
\end{titlepage}
\tableofcontents\newpage


\needspace{12\baselineskip}\section{Lecture at a Glance}

\begin{center}\small\renewcommand{\arraystretch}{1.25}\begin{tabularx}{\textwidth}{@{}>{\bfseries}p{3.3cm}>{\raggedright\arraybackslash}X@{}}\toprule

\textcolor{NavyN}{Lecture}&7 --- Reinforcement Learning Part 2: Policy Gradient Methods\\
\textcolor{NavyN}{Duration}&45 minutes: 5 min opening Q\&A, 30 min teaching, 10 min auto-quiz\\
\textcolor{NavyN}{Slides}&19 (Beamer PDF, 16:9), 10 sections\\
\textcolor{NavyN}{Prerequisite}&Lecture~6; students should have attempted its quiz and assignment.\\
\textcolor{NavyN}{Student materials}&\texttt{Lecture7\_Slides\_Beamer.pdf}; \texttt{Lecture7\_Reference\_Text.pdf}; \texttt{Lecture7\_Quiz.html}; \texttt{Lecture7\_Answers\_Assignment.html} (release after the quiz)\\
\textcolor{NavyN}{Interactive studio}&\texttt{Lecture7\_Policy\_Gradient\_Lab.html} (live demonstrations in Section~\ref{sec:live})\\
\textcolor{NavyN}{Assessment}&Auto-quiz: 15 MCQs, 40~s each, $+1/-0.25$. Assignment: 5 problems, 100 marks, plus a 20-mark bonus problem.\\\bottomrule\end{tabularx}\end{center}

\subsection*{Overview}Lecture~7 moves from value-based to \emph{policy-gradient} RL, motivated by continuous control: a four-drone swarm whose policies output thrust and attitude commands. It derives the policy-gradient theorem and REINFORCE, tames its variance with baselines and advantages, introduces actor--critic (A2C/A3C), and arrives at PPO's clipped surrogate objective. Gaussian policies handle continuous actions; Soft Actor--Critic adds entropy for robustness and sample reuse. The results progress from REINFORCE (61.4\% waypoint success) to PPO (94.3\%) and SAC (96.2\%), and the lecture closes with the sim-to-real gap and domain randomisation (field success 78\% without, 88--91\% with).

\subsection*{Learning outcomes}By the end of the lecture students should be able to:\begin{enumerate}\item Explain why value-based methods struggle with continuous actions.\item State the policy-gradient theorem and the REINFORCE update; explain the score function.\item Prove that a state-dependent baseline leaves the gradient unbiased, and compute advantages.\item Explain actor--critic, A2C/A3C and PPO clipping; compute a clipped objective term.\item Describe Gaussian policies and SAC's entropy objective; plan sim-to-real transfer.\end{enumerate}

\subsection*{Syllabus mapping}\begin{itemize}\item Reinforcement learning for autonomous systems --- \textbf{policy gradient methods} (REINFORCE, actor--critic, PPO, SAC)\end{itemize}

\needspace{12\baselineskip}\section{Session Plan (45 Minutes)}

\begin{center}\begin{tikzpicture}[x=0.3444cm,y=1cm]
\fill[AmbN!40](0,0)rectangle(5,1);\node[font=\scriptsize\bfseries,align=center]at(2.5,0.5){Q\&A\\recap};
\fill[AccN!25](5.00,0)rectangle(7.50,1);\draw[white,thick](5.00,0)--(5.00,1);
\node[font=\tiny,rotate=90,anchor=center]at(6.25,0.5){\parbox{0.9cm}{\centering \S1}};
\fill[NavyN!18](7.50,0)rectangle(11.00,1);\draw[white,thick](7.50,0)--(7.50,1);
\node[font=\tiny,rotate=90,anchor=center]at(9.25,0.5){\parbox{0.9cm}{\centering \S2}};
\fill[AccN!25](11.00,0)rectangle(14.50,1);\draw[white,thick](11.00,0)--(11.00,1);
\node[font=\tiny,rotate=90,anchor=center]at(12.75,0.5){\parbox{0.9cm}{\centering \S3}};
\fill[NavyN!18](14.50,0)rectangle(18.00,1);\draw[white,thick](14.50,0)--(14.50,1);
\node[font=\tiny,rotate=90,anchor=center]at(16.25,0.5){\parbox{0.9cm}{\centering \S4}};
\fill[AccN!25](18.00,0)rectangle(21.50,1);\draw[white,thick](18.00,0)--(18.00,1);
\node[font=\tiny,rotate=90,anchor=center]at(19.75,0.5){\parbox{0.9cm}{\centering \S5}};
\fill[NavyN!18](21.50,0)rectangle(25.00,1);\draw[white,thick](21.50,0)--(21.50,1);
\node[font=\tiny,rotate=90,anchor=center]at(23.25,0.5){\parbox{0.9cm}{\centering \S6}};
\fill[AccN!25](25.00,0)rectangle(26.50,1);\draw[white,thick](25.00,0)--(25.00,1);
\node[font=\tiny,rotate=90,anchor=center]at(25.75,0.5){\parbox{0.9cm}{\centering \S7}};
\fill[NavyN!18](26.50,0)rectangle(30.00,1);\draw[white,thick](26.50,0)--(26.50,1);
\node[font=\tiny,rotate=90,anchor=center]at(28.25,0.5){\parbox{0.9cm}{\centering \S8}};
\fill[AccN!25](30.00,0)rectangle(31.50,1);\draw[white,thick](30.00,0)--(30.00,1);
\node[font=\tiny,rotate=90,anchor=center]at(30.75,0.5){\parbox{0.9cm}{\centering \S9}};
\fill[NavyN!18](31.50,0)rectangle(35.00,1);\draw[white,thick](31.50,0)--(31.50,1);
\node[font=\tiny,rotate=90,anchor=center]at(33.25,0.5){\parbox{0.9cm}{\centering \S10}};
\fill[RedN!30](35,0)rectangle(45,1);\node[font=\scriptsize\bfseries,align=center]at(40,0.5){Auto-quiz\\15 Q $\times$ 40 s};
\draw[NavyN](0,0)rectangle(45,1);
\foreach \tt in {0,5,10,15,20,25,30,35,40,45}{\draw[NavyN](\tt,0)--(\tt,-0.12);\node[font=\tiny,below]at(\tt,-0.12){\tt};}
\node[font=\scriptsize,below]at(22.5,-0.45){minutes};
\end{tikzpicture}\end{center}

\begin{center}\small\renewcommand{\arraystretch}{1.2}\begin{tabularx}{\textwidth}{@{}l l X r@{}}\toprule\textbf{Clock} & \textbf{Slides} & \textbf{Segment} & \textbf{Min}\\\midrule

0:00--5:00 & --- & \textbf{Opening Q\&A and recap} (see prompts below) & 5\\

5:00--7:30 & 2--3 & 1. Why Policy Gradients (2 slides) & 2.5\\

7:30--11:00 & 4--5 & 2. REINFORCE (2 slides) & 3.5\\

11:00--14:30 & 6--7 & 3. Baselines \& Advantage (2 slides) & 3.5\\

14:30--18:00 & 8--9 & 4. Actor-Critic (2 slides) & 3.5\\

18:00--21:30 & 10--11 & 5. PPO (2 slides) & 3.5\\

21:30--25:00 & 12--13 & 6. Continuous Control (2 slides) & 3.5\\

25:00--26:30 & 14 & 7. SAC (1 slide) & 1.5\\

26:30--30:00 & 15--16 & 8. Results (2 slides) & 3.5\\

30:00--31:30 & 17 & 9. Across Disciplines (1 slide) & 1.5\\

31:30--35:00 & 18--19 & 10. Summary (2 slides) & 3.5\\

35:00--45:00 & --- & \textbf{Auto-quiz} (instructions 1 min, 15 questions $\times$ 40 s = 10 min of answering; collect scores) & 10\\\bottomrule\end{tabularx}\end{center}

Slide 1 is the title slide and is shown while students settle. The teaching segment covers 18 content slides in 30 minutes, an average pace of 0.6 slides per minute. 

Minute allocations are proportional to slide counts; shift time towards the sections students find hardest.

\begin{instr}{Opening Q\&A prompts (first 5 minutes)}\begin{enumerate}\item From Lecture~6: why does DQN need a target network?\item What does a $Q$-value of 18.7 mean?\item What does potential-based reward shaping guarantee?\end{enumerate}Take two or three answers per prompt; correct misconceptions briefly and note any topic that needs revisiting.\end{instr}

\needspace{12\baselineskip}\section{Slide-by-Slide Teaching Guide}

Each subsection gives the slide range, the suggested time, what to emphasise, and the slide titles in order. Numbers refer to pages of the Beamer PDF.

\needspace{8\baselineskip}\subsection{1. Why Policy Gradients \normalfont\small\color{GreyN}(slides 2--3, about 2.5 min)}

\textbf{Teaching note.} Why DQN fails for continuous thrust (argmax over a continuum at every step); the intuition of shifting probability toward actions that did well.

\begin{center}\small\renewcommand{\arraystretch}{1.1}\begin{tabularx}{\textwidth}{@{}r X@{}}2 & When DQN Is Not Enough --- Q-Learning Needs a Discrete Action Space\\
3 & Policy Gradient Intuition --- Shifting Probability Toward Good Actions\\\end{tabularx}\end{center}

\needspace{8\baselineskip}\subsection{2. REINFORCE \normalfont\small\color{GreyN}(slides 4--5, about 3.5 min)}

\textbf{Teaching note.} The policy-gradient theorem, the score function $\nabla\log\pi$, REINFORCE with whole-episode returns; the drone set-up and baseline result (61.4\%, never reached 70\%).

\begin{center}\small\renewcommand{\arraystretch}{1.1}\begin{tabularx}{\textwidth}{@{}r X@{}}4 & REINFORCE (Williams 1992) --- Complete Episodes, Then Update\\
5 & Drone Swarm Case Study --- Setup and the REINFORCE Baseline\\\end{tabularx}\end{center}

\needspace{8\baselineskip}\subsection{3. Baselines \& Advantage \normalfont\small\color{GreyN}(slides 6--7, about 3.5 min)}

\textbf{Teaching note.} Why a high return does not mean the action was good; subtracting $b(s)$ keeps the gradient unbiased and halves variance; advantage $A=Q-V$.

\begin{center}\small\renewcommand{\arraystretch}{1.1}\begin{tabularx}{\textwidth}{@{}r X@{}}6 & The Variance Problem --- A High Return Does Not Mean the Action Was Good\\
7 & Baselines in Practice --- Variance Halved With No Bias\\\end{tabularx}\end{center}

\needspace{8\baselineskip}\subsection{4. Actor-Critic \normalfont\small\color{GreyN}(slides 8--9, about 3.5 min)}

\textbf{Teaching note.} Actor and critic; TD advantage $r+\gamma V(s')-V(s)$; A2C (synchronous) versus A3C (asynchronous).

\begin{center}\small\renewcommand{\arraystretch}{1.1}\begin{tabularx}{\textwidth}{@{}r X@{}}8 & Actor--Critic --- Two Networks, One Learning Loop\\
9 & A2C and A3C --- Parallel Workers for Lower Variance\\\end{tabularx}\end{center}

\needspace{8\baselineskip}\subsection{5. PPO \normalfont\small\color{GreyN}(slides 10--11, about 3.5 min)}

\textbf{Teaching note.} Trust-region intuition; the probability ratio and the clipped surrogate; worked: $\varepsilon=0.2$, $r=1.58$, $\hat A=6.8$ gives $\min(10.74,8.16)=8.16$.

\begin{center}\small\renewcommand{\arraystretch}{1.1}\begin{tabularx}{\textwidth}{@{}r X@{}}10 & Proximal Policy Optimisation --- Why Updates Must Stay Close\\
11 & PPO --- The Clipped Surrogate Objective\\\end{tabularx}\end{center}

\needspace{8\baselineskip}\subsection{6. Continuous Control \normalfont\small\color{GreyN}(slides 12--13, about 3.5 min)}

\textbf{Teaching note.} Gaussian policy with learned mean and log-standard-deviation; tanh squashing; what a 6-DoF drone policy controls.

\begin{center}\small\renewcommand{\arraystretch}{1.1}\begin{tabularx}{\textwidth}{@{}r X@{}}12 & Continuous Action Spaces --- the Gaussian Policy\\
13 & 6-DoF Flight Dynamics --- What the Drone Policy Controls\\\end{tabularx}\end{center}

\needspace{8\baselineskip}\subsection{7. SAC \normalfont\small\color{GreyN}(slide 14, about 1.5 min)}

\textbf{Teaching note.} Maximum-entropy objective; robustness under wind (SAC 88\% vs PPO 71\%); off-policy replay.

\begin{center}\small\renewcommand{\arraystretch}{1.1}\begin{tabularx}{\textwidth}{@{}r X@{}}14 & Soft Actor--Critic --- Optimise Reward AND Entropy\\\end{tabularx}\end{center}

\needspace{8\baselineskip}\subsection{8. Results \normalfont\small\color{GreyN}(slides 15--16, about 3.5 min)}

\textbf{Teaching note.} Full results table; choosing a method by sample cost and robustness; domain randomisation over wind, noise, mass and motors.

\begin{center}\small\renewcommand{\arraystretch}{1.1}\begin{tabularx}{\textwidth}{@{}r X@{}}15 & Full Results --- REINFORCE $\to$ PPO $\to$ SAC on the 4-Drone Swarm\\
16 & Choosing the Method --- and Crossing the Sim-to-Real Gap\\\end{tabularx}\end{center}

\needspace{8\baselineskip}\subsection{9. Across Disciplines \normalfont\small\color{GreyN}(slide 17, about 1.5 min)}

\textbf{Teaching note.} Policy learning in neuroscience (motor learning), robotics, aerospace (adaptive flight control) and psychology (habit formation).

\begin{center}\small\renewcommand{\arraystretch}{1.1}\begin{tabularx}{\textwidth}{@{}r X@{}}17 & Policy Gradients \& Continuous Control Across the Disciplines\\\end{tabularx}\end{center}

\needspace{8\baselineskip}\subsection{10. Summary \normalfont\small\color{GreyN}(slides 18--19, about 3.5 min)}

\textbf{Teaching note.} Summary and look-ahead to multi-agent RL.

\begin{center}\small\renewcommand{\arraystretch}{1.1}\begin{tabularx}{\textwidth}{@{}r X@{}}18 & Lecture 7 Summary --- From REINFORCE to SAC for Continuous Drone Control\\
19 & Looking Ahead --- Lecture 8: Multi-Agent RL \& Autonomous Combat Platforms\\\end{tabularx}\end{center}

\needspace{12\baselineskip}\section{Key Messages}Students should leave the room able to state these in their own words:\begin{keypoint}\begin{enumerate}\item Policy gradients optimise the policy directly and handle continuous actions naturally.\item Variance is the enemy: baselines, critics and advantages are all variance-reduction devices.\item PPO's clipping keeps each update close to the previous policy; it is the robust default.\item SAC's entropy term buys robustness and sample reuse.\item Simulation success must be stress-tested; domain randomisation narrows the sim-to-real gap.\end{enumerate}\end{keypoint}

\needspace{12\baselineskip}\section{Facilitation Tips}\begin{itemize}\item Use the ``credit a lucky win to every move'' story to motivate baselines before showing any algebra.\item Compute the PPO clip example live; ask what happens when $\hat A<0$.\item If the class is mathematically strong, do the baseline-unbiasedness proof ($\sum_a\nabla\pi=0$) on the board.\end{itemize}

\needspace{12\baselineskip}\section{Anticipated Questions and Model Answers}

\begin{question}{anticipated question}\textbf{Q:} PPO or SAC for a new project?\par\smallskip\textbf{A:} PPO if simulation is cheap and stability matters; SAC if real interactions are expensive (off-policy reuse) or robustness to disturbance matters.\end{question}

\begin{question}{anticipated question}\textbf{Q:} Why did SAC need more episodes than PPO to reach 70\% if it is more sample-efficient?\par\smallskip\textbf{A:} Episodes are not the scarce resource SAC saves; it reuses each transition many times. Compare on environment steps, compute or wall-clock as appropriate.\end{question}

\begin{question}{anticipated question}\textbf{Q:} Is domain randomisation enough for deployment?\par\smallskip\textbf{A:} No. It narrows the gap; field trials, safety envelopes, monitoring and human authority remain necessary.\end{question}

\needspace{12\baselineskip}\section{Common Misconceptions}\begin{center}\small\renewcommand{\arraystretch}{1.35}\begin{tabularx}{\textwidth}{@{}>{\raggedright\arraybackslash}X >{\raggedright\arraybackslash}X@{}}\toprule\textbf{Misconception} & \textbf{Correction}\\\midrule``Baselines bias the gradient.'' & A state-only baseline has zero expected effect on the gradient.\\
``PPO clips the reward.'' & It clips the probability ratio in the objective.\\
``Higher entropy means random behaviour at deployment.'' & Entropy is a training regulariser; deploy the mean action or a low-temperature policy.\\\bottomrule\end{tabularx}\end{center}

\needspace{12\baselineskip}\section{In-Class Activities and Discussion Prompts}\label{sec:act}\begin{itemize}\item \textbf{Clip table:} give four $(r,\hat A)$ pairs; students compute the clipped term and say whether the gradient flows.\item \textbf{Sim-to-real checklist:} teams list ten parameters they would randomise for a maritime drone.\end{itemize}

\needspace{12\baselineskip}\section{Run It Live: Studio Demonstrations}\label{sec:live}

Project \texttt{Lecture7\_Policy\_Gradient\_Lab.html} (it is a single HTML file and runs offline in any browser; every result updates as an input changes). Each demonstration below follows the same rhythm: \textbf{set} the inputs, \textbf{ask} the class to commit to a prediction (a show of hands or a one-word answer), then \textbf{reveal} the result and \textbf{land} the point. The numbers are exactly what the page shows for the settings given; simulations use fixed random seeds, so they repeat every time. The 45-minute plan has little slack: choose the one demonstration marked \emph{first choice} for class, and use the others in a tutorial, a second session, or as homework with the drill cards at the bottom of the page.

\begin{live}{Demo 1: Why a baseline matters (about 3 min, with \S3 Baselines \& Advantage) --- \emph{first choice}}\textbf{Set.} Demo~1: read the gradient-variance and success rows as the reward offset is set to 0, 10 and 50.\par\smallskip\textbf{Ask.} ``Adding 50 to every reward does not change which route is best. Does it change learning?''\par\smallskip\textbf{Reveal.} Gradient variance without/with a baseline: offset 0: 2.2 / 0.72; offset 10: 88.9 / 0.72; offset 50: \textbf{1,768.9 / 0.72}. Runs (of 40) ending on the best route: 39/40, 21/40, \textbf{14/40} without a baseline; 40/40 with one every time.\par\smallskip\textbf{Land the point.} Subtracting a baseline leaves the expected gradient unchanged and removes most of the noise; without it, early lucky choices lock in.\end{live}

\begin{live}{Demo 2: Entropy collapse in a Gaussian policy (about 3 min, with \S6 Continuous Control or \S7 SAC)}\textbf{Set.} Demo~2: entropy bonus $\beta=0$; press \textsf{Run 50 updates} four times. Then set $\beta=0.5$, press \textsf{Reset policy}, and run six times.\par\smallskip\textbf{Ask.} ``The aim converges on 3. What happens to the spread?''\par\smallskip\textbf{Reveal.} $\beta=0$: $\mu=2.979$, $\sigma=0.161$ and still shrinking. $\beta=0.5$: $\mu=3.059$, $\sigma=$ \textbf{0.534}, near the optimum $\sqrt{\beta/2}=0.5$.\par\smallskip\textbf{Land the point.} Without an entropy term the policy stops exploring. The bonus buys exploration at a known cost ($\sigma^2$ of expected reward here).\end{live}

\begin{live}{Demo 3: When does a PPO sample stop pushing? (about 2 min, with \S5 PPO)}\textbf{Set.} Demo~3: $\varepsilon=0.2$. Try $(A,\rho)=(+1,1.35)$, $(-1,0.6)$ and $(-1,1.4)$.\par\smallskip\textbf{Ask.} For each case: ``Is the gradient zero?''\par\smallskip\textbf{Reveal.} $(+1,1.35)$: objective 1.200, gradient \textbf{0} (clipped). $(-1,0.6)$: $-0.800$, gradient \textbf{0}. $(-1,1.4)$: $-1.400$, gradient \textbf{$-1$}.\par\smallskip\textbf{Land the point.} Clipping only removes the incentive to move further in the direction the advantage already favours; mistakes in the wrong direction are still corrected.\end{live}

\needspace{12\baselineskip}\section{Consistency Notes on the Source Material}These notes record where the original PowerPoint deck contained an inconsistency and how the current slides, reference text, quiz and assignment handle it. Students who have seen the old deck may raise them.\begin{itemize}\item The source deck's comparison table lists 6,200 ``episodes to 70\%'' for REINFORCE, but REINFORCE peaked at 61.4\% and never reached 70\%; 6,200 is the number of episodes to reach 50\%. The new slides and reference text say so.\end{itemize}

\needspace{12\baselineskip}\section{Quiz Administration and Answer Key}

\begin{sayit}``You will now take a 15-question quiz. Each question has four options and exactly one correct answer. You have 40 seconds per question; the quiz advances automatically when time runs out. A correct answer scores $+1$, a wrong answer $-0.25$, and an unanswered question 0. The maximum is 15 and the minimum is $-3.75$. Because of the negative marking, guessing blindly among four options gains nothing on average --- answer when you can eliminate options. Your score and a full review appear at the end.''\end{sayit}

\textbf{Delivery.} Open \texttt{Lecture7\_Quiz.html} on each student device (or project it and collect answers on paper). The page enforces the timer and computes the score with negative marking. Release \texttt{Lecture7\_Answers\_Assignment.html} only after all students have finished.

\subsection*{Answer key}\begin{center}\small\renewcommand{\arraystretch}{1.3}\begin{tabular}{|c|c|c|c|c|c|c|c|c|c|c|c|c|c|c|}\hline \textbf{1} & \textbf{2} & \textbf{3} & \textbf{4} & \textbf{5} & \textbf{6} & \textbf{7} & \textbf{8} & \textbf{9} & \textbf{10} & \textbf{11} & \textbf{12} & \textbf{13} & \textbf{14} & \textbf{15}\\\hline A & C & C & D & C & D & B & A & B & D & D & A & C & B & B\\\hline\end{tabular}\end{center}

\textbf{Interpreting scores.} 12 or more: secure command of the lecture. 9--11.75: good, review missed items. Below 9: revisit the reference text before the assignment. If more than half the class misses a question, revisit that point in the next lecture's opening Q\&A.

\subsection*{Questions, options and rationale}

{\small

\noindent\textbf{Q1.} Policy-gradient methods are preferred over DQN for drone thrust control because:\begin{itemize}[leftmargin=2.2em,itemsep=0pt,topsep=1pt]\item[\textbf{A.}] \textbf{Continuous actions make argmax over Q an optimisation at every step; a policy can output a distribution directly}~{\color{GrnN}$\checkmark$}\item[B.] They need no reward\item[C.] They are always faster\item[D.] They need no exploration\end{itemize}\noindent{\color{GreyN}\textit{Rationale:}} A Gaussian policy outputs mean and standard deviation; actions are sampled.\par\medskip

\noindent\textbf{Q2.} The policy gradient is $\mathbb E[\nabla\log\pi_\theta(a|s)\,G_t]$. The term $\nabla\log\pi$ is called the:\begin{itemize}[leftmargin=2.2em,itemsep=0pt,topsep=1pt]\item[A.] Advantage\item[B.] Entropy\item[\textbf{C.}] \textbf{Score function}~{\color{GrnN}$\checkmark$}\item[D.] TD error\end{itemize}\noindent{\color{GreyN}\textit{Rationale:}} It points in the direction that makes the chosen action more likely.\par\medskip

\noindent\textbf{Q3.} REINFORCE's main weakness is:\begin{itemize}[leftmargin=2.2em,itemsep=0pt,topsep=1pt]\item[A.] Bias\item[B.] It needs a model\item[\textbf{C.}] \textbf{High variance from using whole-episode returns}~{\color{GrnN}$\checkmark$}\item[D.] It cannot handle discrete actions\end{itemize}\noindent{\color{GreyN}\textit{Rationale:}} Return spread was 84.6 on the drone task.\par\medskip

\noindent\textbf{Q4.} Subtracting a baseline b(s) from the return:\begin{itemize}[leftmargin=2.2em,itemsep=0pt,topsep=1pt]\item[A.] Biases the gradient\item[B.] Removes exploration\item[C.] Changes the optimal policy\item[\textbf{D.}] \textbf{Leaves the expected gradient unchanged and reduces variance}~{\color{GrnN}$\checkmark$}\end{itemize}\noindent{\color{GreyN}\textit{Rationale:}} $\mathbb E[\nabla\log\pi\,b(s)]=0$.\par\medskip

\noindent\textbf{Q5.} A positive advantage A(s,a) means:\begin{itemize}[leftmargin=2.2em,itemsep=0pt,topsep=1pt]\item[A.] The critic is wrong\item[B.] The state is good\item[\textbf{C.}] \textbf{The action was better than average from that state; increase its probability}~{\color{GrnN}$\checkmark$}\item[D.] Decrease the action's probability\end{itemize}\noindent{\color{GreyN}\textit{Rationale:}} $A=Q-V$.\par\medskip

\noindent\textbf{Q6.} In actor-critic, the one-step advantage estimate is:\begin{itemize}[leftmargin=2.2em,itemsep=0pt,topsep=1pt]\item[A.] $G_t$\item[B.] $\log\pi$\item[C.] $Q(s,a)$\item[\textbf{D.}] \textbf{$r+\gamma V(s')-V(s)$}~{\color{GrnN}$\checkmark$}\end{itemize}\noindent{\color{GreyN}\textit{Rationale:}} The TD error: lower variance, some bias from an imperfect critic.\par\medskip

\noindent\textbf{Q7.} A2C differs from A3C in that A2C:\begin{itemize}[leftmargin=2.2em,itemsep=0pt,topsep=1pt]\item[A.] Uses no critic\item[\textbf{B.}] \textbf{Synchronises updates from parallel workers (no stale gradients)}~{\color{GrnN}$\checkmark$}\item[C.] Is off-policy\item[D.] Uses a replay buffer\end{itemize}\noindent{\color{GreyN}\textit{Rationale:}} A3C updates asynchronously and can suffer from stale gradients.\par\medskip

\noindent\textbf{Q8.} PPO clips the probability ratio to $[1-\varepsilon,1+\varepsilon]$ in order to:\begin{itemize}[leftmargin=2.2em,itemsep=0pt,topsep=1pt]\item[\textbf{A.}] \textbf{Prevent destructively large policy updates}~{\color{GrnN}$\checkmark$}\item[B.] Increase exploration\item[C.] Remove the critic\item[D.] Scale rewards\end{itemize}\noindent{\color{GreyN}\textit{Rationale:}} First-order trust-region behaviour without TRPO's second-order cost.\par\medskip

\noindent\textbf{Q9.} With $\varepsilon=0.2$, $r_t=1.58$, $\hat A=6.8$, the clipped objective term is:\begin{itemize}[leftmargin=2.2em,itemsep=0pt,topsep=1pt]\item[A.] 10.76\item[\textbf{B.}] \textbf{8.16}~{\color{GrnN}$\checkmark$}\item[C.] 6.8\item[D.] 1.58\end{itemize}\noindent{\color{GreyN}\textit{Rationale:}} $\min(1.58\times6.8,1.2\times6.8)=8.16$.\par\medskip

\noindent\textbf{Q10.} In a Gaussian policy, the network typically outputs:\begin{itemize}[leftmargin=2.2em,itemsep=0pt,topsep=1pt]\item[A.] Q-values\item[B.] Rewards\item[C.] A single deterministic action\item[\textbf{D.}] \textbf{Mean and log-standard-deviation of the action distribution}~{\color{GrnN}$\checkmark$}\end{itemize}\noindent{\color{GreyN}\textit{Rationale:}} Log-sigma keeps sigma positive and stabilises training.\par\medskip

\noindent\textbf{Q11.} SAC's objective adds which term to the reward?\begin{itemize}[leftmargin=2.2em,itemsep=0pt,topsep=1pt]\item[A.] A KL penalty\item[B.] A clip\item[C.] A baseline\item[\textbf{D.}] \textbf{$\alpha H(\pi(\cdot|s))$, the policy entropy}~{\color{GrnN}$\checkmark$}\end{itemize}\noindent{\color{GreyN}\textit{Rationale:}} Maximum-entropy RL: explore, stay robust, stay unpredictable.\par\medskip

\noindent\textbf{Q12.} Under wind perturbation PPO fell to 71\% while SAC held 88\%. The best explanation is:\begin{itemize}[leftmargin=2.2em,itemsep=0pt,topsep=1pt]\item[\textbf{A.}] \textbf{SAC's entropy term learns many ways to succeed, so it is robust to disturbances}~{\color{GrnN}$\checkmark$}\item[B.] SAC used more data\item[C.] PPO is broken\item[D.] Wind affects PPO only\end{itemize}\noindent{\color{GreyN}\textit{Rationale:}} Stochastic, diverse behaviour generalises better to perturbations.\par\medskip

\noindent\textbf{Q13.} SAC is off-policy, which means:\begin{itemize}[leftmargin=2.2em,itemsep=0pt,topsep=1pt]\item[A.] It needs complete episodes\item[B.] It never explores\item[\textbf{C.}] \textbf{It learns from a replay buffer of past experience, reusing data}~{\color{GrnN}$\checkmark$}\item[D.] It has no critic\end{itemize}\noindent{\color{GreyN}\textit{Rationale:}} Hence high sample efficiency per environment interaction.\par\medskip

\noindent\textbf{Q14.} Drone success fell from 94\% in simulation to 78\% in the field. The recommended remedy is:\begin{itemize}[leftmargin=2.2em,itemsep=0pt,topsep=1pt]\item[A.] Train longer in the same simulator\item[\textbf{B.}] \textbf{Domain randomisation over wind, noise, mass and motor variations}~{\color{GrnN}$\checkmark$}\item[C.] Use REINFORCE\item[D.] Lower gamma\end{itemize}\noindent{\color{GreyN}\textit{Rationale:}} With randomisation real-world success rose to 88\% (PPO) and 91\% (SAC).\par\medskip

\noindent\textbf{Q15.} The SAC drones kept about 32.4 m spacing without being told. This illustrates:\begin{itemize}[leftmargin=2.2em,itemsep=0pt,topsep=1pt]\item[A.] Overfitting\item[\textbf{B.}] \textbf{Emergent coordinated behaviour from reward plus entropy}~{\color{GrnN}$\checkmark$}\item[C.] A bug\item[D.] Reward hacking\end{itemize}\noindent{\color{GreyN}\textit{Rationale:}} Spacing between the 10 m penalty and 50 m bonus emerged naturally.\par\medskip

}

\needspace{12\baselineskip}\section{Assignment: Marking Scheme and Model Solutions}

\textbf{Instructions.} Five problems, 100 marks, plus a 20-mark bonus problem (Problem 6). Reference: Lecture 7 reference text.

\begin{center}\small\renewcommand{\arraystretch}{1.25}\begin{tabularx}{0.85\textwidth}{@{}l X r@{}}\toprule\textbf{\#} & \textbf{Problem} & \textbf{Marks}\\\midrule1 & REINFORCE by hand & 20\\
2 & The baseline proof and advantages & 20\\
3 & PPO clipping & 20\\
4 & Gaussian policy update & 20\\
5 & Choosing and validating a controller & 20\\
6 & Problem 6 (bonus): Policy-gradient updates by hand & 20\\\midrule & \textbf{Total} & \textbf{100} + 20 bonus\\\bottomrule\end{tabularx}\end{center}

\begin{instr}{General marking guidance}\begin{itemize}
\item Marks for each part are shown in square brackets in the problem statement.
\item Award method marks for correct working even when an arithmetic slip carries through; deduct once for the slip, not again for its consequences.
\item For interpretive parts, full marks need (i) the correct conclusion, (ii) the reason, and (iii) an operational implication or caveat.
\item Answers that state a number without uncertainty, baseline or caveat where one is required lose a mark per part.
\item Suggested deadline: one week after the lecture; return marked work before the next quiz on the same topic.
\end{itemize}\end{instr}

\needspace{8\baselineskip}\subsection{Problem 1: REINFORCE by hand \normalfont\small\color{GreyN}(20 marks)}

\begin{prob}A softmax policy over \{left, right\} has preferences $h=(0,0)$; $\nabla_h\log\pi(a)=\mathbf 1_a-\pi$.

(a) The agent takes \textit{right} and receives return $G=4$; with $\alpha=0.1$ compute the new preferences and $\pi(\text{right})$. [8]\newline (b) Repeat from the original policy with a baseline $b=3$. [6]\newline (c) Explain why the baseline version is preferable when returns range between 0 and 8. [6]\end{prob}

\begin{soln}\textbf{(a)} $\nabla\log\pi(\text{right})=(-0.5,0.5)$. $h\leftarrow0.1\times4\times(-0.5,0.5)=(-0.2,0.2)$; $\pi(\text{right})=e^{0.2}/(e^{0.2}+e^{-0.2})=\mathbf{0.60}$.

\textbf{(b)} Effective weight $G-b=1$: $h=(-0.05,0.05)$; $\pi(\text{right})=e^{0.05}/(e^{0.05}+e^{-0.05})=\mathbf{0.525}$.

\textbf{(c)} Without a baseline every action is reinforced (all returns are positive) and only the relative size of updates carries information, so variance is high. With a baseline near the mean return, below-average outcomes push probability down and above-average ones push it up; updates are smaller and better-signed, and the expected gradient is unchanged.\end{soln}

\needspace{8\baselineskip}\subsection{Problem 2: The baseline proof and advantages \normalfont\small\color{GreyN}(20 marks)}

\begin{prob}(a) Prove $\mathbb E_{a\sim\pi}[\nabla_\theta\log\pi_\theta(a|s)\,b(s)]=0$. [8]\newline (b) Why must the baseline not depend on the action? [4]\newline (c) A critic gives $V(s)=10$, $V(s')=12$; reward $r=1$, $\gamma=0.9$. Compute the one-step advantage, and explain the sign. [8]\end{prob}

\begin{soln}\textbf{(a)} $\sum_a\pi(a|s)\frac{\nabla\pi(a|s)}{\pi(a|s)}b(s)=b(s)\nabla\sum_a\pi(a|s)=b(s)\nabla1=0$.

\textbf{(b)} If $b$ depended on $a$ it could not be taken outside the sum, the identity would fail, and the gradient would be biased.

\textbf{(c)} $\hat A=1+0.9\times12-10=\mathbf{1.8}$. The transition led to a state worth more than expected, so the action was better than average and its probability is increased.\end{soln}

\needspace{8\baselineskip}\subsection{Problem 3: PPO clipping \normalfont\small\color{GreyN}(20 marks)}

\begin{prob}$\varepsilon=0.2$. Compute the clipped-objective term and state whether its gradient with respect to $r$ is zero, for: (i) $r=1.1,\hat A=5$; (ii) $r=1.5,\hat A=5$; (iii) $r=0.9,\hat A=-3$; (iv) $r=0.6,\hat A=-3$; (v) $r=1.5,\hat A=-3$. [15] Explain case (v). [5]\end{prob}

\begin{soln}(i) $\min(5.5,5.5)=5.5$, gradient non-zero. (ii) $\min(7.5,6.0)=\mathbf{6.0}$, clipped, zero gradient. (iii) $\min(-2.7,-2.7)=-2.7$, non-zero. (iv) $\min(-1.8,-2.4)=\mathbf{-2.4}$, clipped, zero gradient. (v) $\min(-4.5,-3.6)=\mathbf{-4.5}$, \textbf{not} clipped.

Case (v): a bad action has become much \textit{more} likely. The min keeps the unclipped, more pessimistic term, so the gradient still pushes its probability down. PPO clips only in the direction that would over-reward a change, never in the direction that corrects a mistake.\end{soln}

\needspace{8\baselineskip}\subsection{Problem 4: Gaussian policy update \normalfont\small\color{GreyN}(20 marks)}

\begin{prob}A 1-D thrust policy has $\mu=2.0$, $\sigma=0.5$.

(a) Sampled $a=2.4$, advantage $+3$, step size 0.01: update $\mu$. [6]\newline (b) Sampled $a=1.5$, advantage $-2$: update $\mu$ from 2.0. [6]\newline (c) Derive $\partial\log\pi/\partial\sigma$ and explain how $\sigma$ changes when an action far from the mean has positive advantage. [8]\end{prob}

\begin{soln}\textbf{(a)} $\partial\log\pi/\partial\mu=(a-\mu)/\sigma^2=1.6$; $\mu=2.0+0.01\times3\times1.6=\mathbf{2.048}$.

\textbf{(b)} $(1.5-2.0)/0.25=-2$; times $\hat A=-2$ gives $+4$; $\mu=2.0+0.04=\mathbf{2.04}$ --- away from the bad action.

\textbf{(c)} $\log\pi=-\frac{(a-\mu)^2}{2\sigma^2}-\log\sigma+\text{const}$, so $\partial/\partial\sigma=\frac{(a-\mu)^2}{\sigma^3}-\frac1\sigma$. It is positive when $|a-\mu|>\sigma$. A distant action with positive advantage therefore increases $\sigma$ (explore more widely); a near-mean success decreases it (act more precisely).\end{soln}

\needspace{8\baselineskip}\subsection{Problem 5: Choosing and validating a controller \normalfont\small\color{GreyN}(20 marks)}

\begin{prob}You must field an autonomous waypoint controller for a 4-drone ISR swarm.

(a) Compare PPO and SAC on five criteria and recommend one. [8]\newline (b) Design a domain-randomisation scheme and a field-validation protocol. [8]\newline (c) What safety architecture would you place around the learned policy? [4]\end{prob}

\begin{soln}\textbf{(a)} Criteria: sample efficiency (SAC off-policy reuse vs PPO on-policy); robustness (SAC 88\% vs PPO 71\% in wind); stability and tuning effort (PPO simpler); predictability for certification (PPO near-deterministic; SAC must anneal its stochasticity); adversarial unpredictability (SAC better). For a contested ISR mission where robustness matters and simulation is cheap, SAC with entropy annealed at deployment, keeping PPO as a baseline.

\textbf{(b)} Randomise wind (0--8 m/s, gusts), sensor and GPS noise ($\pm2$ m), mass ($\pm15$\%), motor coefficients ($\pm10$\%), latency (0--50 ms) and battery sag. Validate in hardware-in-the-loop, then tethered and range trials across weather, reporting success rate with confidence intervals over many sorties and a pre-set acceptance threshold.

\textbf{(c)} A certified PID inner loop; geofence and minimum-separation enforcement outside the learned policy; a fallback controller and return-to-home on anomaly; human override and full logging. The learned policy proposes waypoints and thrust targets within a verified safety envelope.\end{soln}

\needspace{8\baselineskip}\subsection{Problem 6 (bonus): Policy-gradient updates by hand \normalfont\small\color{GreyN}(20 marks)}

\begin{prob}(a) A climb/hold policy has preferences $(0.2,-0.2)$. The agent climbed and the episode return was $G=10$; the baseline is $b=6$. With $\alpha=0.05$, update the preferences and the probabilities. Repeat without the baseline and comment. [6]\newline (b) A Gaussian thrust policy has $\mu=0.40$, $\sigma=0.10$; it sampled $a=0.47$ and the advantage was $+3$. Compute the gradients with respect to $\mu$ and $\ln\sigma$, and the new $\mu$ with $\alpha=0.001$. [5]\newline (c) PPO with $\varepsilon=0.2$: three samples have $(r_t,\hat A_t)=(1.35,+2.0)$, $(0.70,-1.5)$ and $(1.10,-0.8)$. Compute each clipped term and their mean. Which samples contribute no gradient? [6]\newline (d) Rewards $1,0,5$ over three steps that end the episode; value estimates $V=2.0,2.5,3.0$; $\gamma=0.99$, $\lambda=0.95$. Compute the GAE advantages. [3]\end{prob}

\begin{soln}\textbf{(a)} $\pi=(0.599,0.401)$; score $(1-0.599,\;-0.401)=(0.401,-0.401)$. With advantage $G-b=4$: preferences change by $0.05\times4\times(0.401,-0.401)$ to $(0.280,-0.280)$, giving $\pi=(\mathbf{0.637},0.363)$. Without the baseline (multiplier 10): $\pi(\text{climb})=0.690$. The direction is the same, but the step is 2.5 times larger. Because every return is positive, an unbaselined update would push up whatever action was taken; the baseline credits climbing only for doing better than usual, which cuts variance.

\textbf{(b)} $\partial\ln\pi/\partial\mu=(0.47-0.40)/0.01=7$; times $A=3$: $\mathbf{21}$, so $\mu\to0.40+0.001\times21=\mathbf{0.421}$. $\partial\ln\pi/\partial\ln\sigma=(0.07^2/0.01-1)=-0.51$; times 3: $\mathbf{-1.53}$, so $\sigma$ shrinks --- a good action close to the mean narrows the distribution.

\textbf{(c)} Sample 1: $\min(1.35\times2,\;1.2\times2)=\mathbf{2.4}$ (clipped). Sample 2: $\min(0.70\times-1.5,\;0.8\times-1.5)=\min(-1.05,-1.2)=\mathbf{-1.2}$ (clipped). Sample 3: $r$ is inside $[0.8,1.2]$, so the term is $\mathbf{-0.88}$. Mean $=\mathbf{0.107}$. Samples 1 and 2 take the clipped constant, so they contribute no gradient: the policy has already moved as far as PPO allows in those directions. Only sample 3 drives the update.

\textbf{(d)} TD errors: $\delta_0=1+0.99(2.5)-2=1.475$; $\delta_1=0+0.99(3)-2.5=0.470$; $\delta_2=5+0-3=2.0$. With $\gamma\lambda=0.9405$: $\hat A_2=\mathbf{2.00}$, $\hat A_1=0.470+0.9405(2.0)=\mathbf{2.35}$, $\hat A_0=1.475+0.9405(2.351)=\mathbf{3.69}$.\end{soln}

\end{document}